\section{安全治理与审查流程}
\subsection{审计链设计}
OpenHands 生成的每一次命令、文件 diff、测试结果都被送入审计链，由审计服务写入事件中间件。这让团队可以在事后轻松查看“谁让 Agent 执行了哪条命令”，并与安全策略对照。

\begin{table}[H]
\centering
\begin{tabular}{p{0.25\textwidth} p{0.4\textwidth} p{0.33\textwidth}}
\toprule
审计阶段 & 内容 & 输出 \\
\midrule
事件捕捉 & 命令、浏览、文件操作 & Log ID + 时间戳 \\
\midrule
策略對齊 & 检查是否触发敏感资源 & Warning / Approval request \\
\midrule
报告归档 & 生成邮件/issue & 可复现的审核记录 \\
\bottomrule
\end{tabular}
\caption{安全审计流程中的三个阶段。}
\end{table}

\begin{importantbox}{审计要素}
记录要包含上下文（Issue、diff、测试），便于人类快速判断风险。
\end{importantbox}

\subsection{审批与权限}
当 Agent 计划执行敏感命令（比如数据库迁移、云部署）时，审批模块会自动触发。审批决策包含触发命令、调用环境、预期结果、相关 diff 以及上一次的测试状态，审批者只需确认或拒绝。

\begin{warningbox}{审批的节奏}
不要过度审批以致阻塞迭代；而在高风险场景必须设置 multi-party approval，防止单点失效。
\end{warningbox}

\subsection{本章小结}
\begin{itemize}
    \item OpenHands 的审计链把 Agent 行为录下来，用于安全复盘。
    \item 合理设置审批与权限可以在不降低效率的情况下强化安全。
\end{itemize}

